\documentclass[12pt]{ximera}
\typeout{Start loading xmPreamble.tex}%

% Add here extra macro's that are loaded automatically by all documents of claas 'ximera' or 'xourse' in this repo

%%
%%  Example:
%%
% \newcommand{\R}{\mathbb{R}}
\usepackage{tikz}
\usetikzlibrary{positioning, arrows.meta}
\usepackage{multicol}
\usepackage{enumitem}
\usepackage{caption}
\usepackage{amsmath}
\usepackage{amsthm}

%% NOTE: do NOT add \usepackage to individual activity .tex files.
%% When a xourse inlines an activity it gobbles \documentclass and \input, but a bare
%% \usepackage still executes -- after \begin{document} -- and errors out.
%% All shared packages and macros belong here.
%%
%% ximera.cls already loads: amsmath, amssymb, amsthm, cancel, enumitem, graphicx,
%% hyperref, listings, pgfplots, tikz, url, xcolor. No need to re-load those.

\title{In-Class: Linear Functions and Models}
\begin{document}
\begin{abstract}
Building and interpreting linear models in context: slope-intercept form, what the
slope and intercept mean (and what units they carry), electricity billing, comparing
the cost of a conventional and a hybrid car, and reading a linear regression line
relating ulna length to height.
\end{abstract}
\maketitle

\section*{Lines in slope-intercept form}

\[
y = \underbrace{m}_{\text{slope}} x + \underbrace{b}_{y\text{-intercept}}
\]

\[
\text{slope} = \frac{\text{rise}}{\text{run}}
= \frac{\text{change in } y\text{-values}}{\text{change in } x\text{-values}}
\]

If we increase the $x$-value by $1$, then the $y$-value increases by $m$.

The \textbf{$y$-intercept} is, graphically, the $y$-coordinate of the point where the
line intersects the $y$-axis. The \textbf{$x$-intercept} is the $x$-coordinate of the
point where the line intersects the $x$-axis.

We're going to revisit some problems from your pre-class work, and the goal today is
to understand and interpret these functions \emph{in the context of the problems}.

\subsection*{A wedding budget}

\begin{problem}
Your friends are trying to get an idea of how many people they can invite to their
wedding. The venue they're looking at costs about \$2,000 to rent, and the event
coordinator suggested they budget about \$100 per attendee.

\begin{prompt}
\textbf{(a)} From your pre-class work, record your function for this problem. Let $x$
be the number of attendees and $y$ the total cost in dollars.
\[
y = \answer{100} x + \answer{2000}
\]

\begin{solution}
With units made explicit:
\[
\$y = \frac{\$100}{\text{attendee}} \cdot x \text{ attendees} + \$2000.
\]
\end{solution}
\end{prompt}

\begin{prompt}
\textbf{(b)} What is the ``$b$-value'' in your equation? What does it mean in this
context?

$b = \$\answer{2000}$

\begin{freeResponse}
\end{freeResponse}

\begin{solution}
The $b$-value is the $y$-intercept. Algebraically it is the $y$-value when $x = 0$;
graphically it is where the line crosses the $y$-axis. In this context: when the
number of attendees is $0$, the cost of the wedding is still \$2000 --- the venue
rental.
\end{solution}
\end{prompt}

\begin{prompt}
\textbf{(c)} What is the ``$m$-value'' in your equation?
$m = \$\answer{100}$ per attendee.

\textbf{i.} What are the units of the $m$-value? Explain why it makes sense that the
$m$-value has these units, based on the units of $x$ and $y$.
\begin{freeResponse}
\end{freeResponse}

\textbf{ii.} What does this $m$-value mean in this context?
\begin{freeResponse}
\end{freeResponse}

\begin{solution}
Since $\text{slope} = \frac{\text{change in } y}{\text{change in } x}$, the units of
$m$ are always $\frac{\text{units of } y}{\text{units of } x}$ --- here, dollars per
attendee.

Mathematically, increasing the input $x$ by $1$ increases the output $y$ by $100$. In
context: each additional attendee adds \$100 to the cost of the wedding.
\end{solution}
\end{prompt}
\end{problem}

\subsection*{Electricity billing}

\begin{problem}
Below is some information for how electricity usage is billed in Columbus. See
\href{https://www.columbus.gov/}{columbus.gov} for the Schedule A and Schedule A-1
rate information.

% TODO: the rate-schedule image did not survive the PDF conversion (the source file
% was a broken-image placeholder). Re-capture the Schedule A / A-1 table from the
% original worksheet and place it in xmPictures/linearModels/.

\begin{prompt}
\textbf{(a)} In your pre-class work you determined how much it would cost to use
700 KWH using both Schedule A and Schedule A-1. Record that information here.
\begin{freeResponse}
\end{freeResponse}
\end{prompt}

\begin{prompt}
\textbf{(b)} Effectively, about how much does the 700th KWH cost? Why?
\begin{freeResponse}
\end{freeResponse}
\end{prompt}

\begin{prompt}
\textbf{(c)} Explain why, based on this information, there cannot be a monthly bill
for \$65.
\begin{freeResponse}
\end{freeResponse}
\end{prompt}
\end{problem}

\subsection*{Comparing two cars}

\begin{problem}
According to \href{http://www.truecar.com}{TrueCar}, a 2018 Toyota Camry
(conventional, 29/41 MPG city/hwy) sells for an average of \$22,030, and a Camry
Hybrid (51/53 MPG city/hwy) sells for an average of \$26,247. Let's estimate gas at
about \$2.80 per gallon.

\begin{prompt}
\textbf{(a)} From your pre-class work, record your functions for driving each car $x$
miles. Average the city and highway MPG for each vehicle.

Conventional: $y = \answer[tolerance=0.005]{0.08} x + \answer{22030}$

Hybrid: $y = \answer[tolerance=0.005]{0.054} x + \answer{26247}$

\begin{hint}
Think about units. $x$ is in miles and $y$ is in dollars, so $m$ must have units of
$\frac{\$}{\text{mile}}$. You are given $\frac{\text{miles}}{\text{gallon}}$ and
$\frac{\$}{\text{gallon}}$ --- how do you combine them to get $\frac{\$}{\text{mile}}$?
\end{hint}

\begin{solution}
Conventional: average the MPG, $\frac{29 + 41}{2} = 35$ mi/gal. Then
\[
\frac{\$2.80}{1 \cancel{\text{gal}}} \cdot \frac{1 \cancel{\text{gal}}}{35 \text{ miles}}
= \frac{\$0.08}{\text{mile}},
\]
so $m = \$0.08$ per mile and
$\$y = \frac{\$0.08}{1 \text{ mile}} \cdot x \text{ miles} + \$22{,}030$.

Hybrid: average $\frac{51 + 53}{2} = 52$ mi/gal, so
$\frac{\$2.80}{52 \text{ miles}} \approx \frac{\$0.054}{\text{mile}}$, giving
$y \approx 0.054x + 26{,}247$.
\end{solution}
\end{prompt}

\begin{prompt}
\textbf{(b)} What are the ``$b$-values'' in these expressions? What do they represent
in this context?
\begin{freeResponse}
\end{freeResponse}

\begin{solution}
Conventional: $b = \$22{,}030$; hybrid: $b = \$26{,}247$. In context, when the miles
driven is $0$ the cost is just the purchase price of the car.
\end{solution}
\end{prompt}

\begin{prompt}
\textbf{(c)} What are the ``$m$-values'' in these expressions --- write them with
their units. What do they mean in this context?
\begin{freeResponse}
\end{freeResponse}

\begin{solution}
Conventional: $m = \frac{\$0.08}{\text{mile}}$; hybrid:
$m \approx \frac{\$0.054}{\text{mile}}$. Each additional mile driven costs about
8 cents in the conventional Camry and about 5.4 cents in the hybrid.
\end{solution}
\end{prompt}

\begin{prompt}
\textbf{(d)} What other factors could we be considering when comparing the ``costs''
of these two vehicles?
\begin{freeResponse}
\end{freeResponse}
\end{prompt}
\end{problem}

\subsection*{Ulna length and height}

\begin{problem}
In your pre-class work you measured your height and your ulna to the nearest 0.5 cm
and entered that data into the shared spreadsheet linked in Carmen. Your instructor
will use this data to create a scatterplot and a linear regression line. Using the
class data, answer the following.

\begin{prompt}
\textbf{(a)} Record the linear regression line here. (A sample line from a previous
class: $y = 1.9x + 118$. Use your own class's line if it differs.)
\begin{freeResponse}
\end{freeResponse}
\end{prompt}

\begin{prompt}
\textbf{(b)} What is the value of the slope? What does it mean in terms of this
context?
\begin{freeResponse}
\end{freeResponse}

\begin{solution}
The slope is $\frac{1.9 \text{ cm of height}}{1 \text{ cm of ulna}}$. According to the
model, for every increase of 1 cm in ulna length, the model predicts an increase of
1.9 cm in a person's height.
\end{solution}
\end{prompt}

\begin{prompt}
\textbf{(c)} What is the $y$-intercept value? What does it mean in this context? Does
this value make sense in the real world --- why or why not?
\begin{freeResponse}
\end{freeResponse}

\begin{solution}
The $y$-intercept is 118 cm \emph{of height}. In context: when the ulna length is
0 cm, the model predicts the person is 118 cm tall.

Not really meaningful in the real world. An ulna length of 0 cm does not describe
anyone, so the model says nothing useful there --- it is far outside the range of the
data used to build it.
\end{solution}
\end{prompt}

\begin{prompt}
\textbf{(d)} Suppose someone's ulna bone is 24 cm long. According to the model
$y = 1.9x + 118$, how tall is this person likely to be?
$\answer[tolerance=0.05]{163.6}$ cm.

\begin{solution}
$y = 1.9(24) + 118 = 45.6 + 118 = 163.6$ cm.
\end{solution}
\end{prompt}

\begin{prompt}
\textbf{(e)} What is an acceptable amount for your model's height to be off by? Is
off by 5 cm acceptable? Off by 10 cm? By 1 meter? Explain your decision.
\begin{freeResponse}
\end{freeResponse}
\end{prompt}

\begin{prompt}
\textbf{(f)} Suppose someone is 170 cm tall. About how long is their ulna bone likely
to be, according to the model?
$\answer[tolerance=0.05]{27.37}$ cm.

\begin{solution}
\[
170 = 1.9x + 118 \implies 52 = 1.9x \implies x = \frac{52}{1.9} \approx 27.37 \text{ cm}.
\]
Note that here we are given the output and solving for the input --- the reverse of
part (d).
\end{solution}
\end{prompt}
\end{problem}

\subsection*{A scientific study}

Researchers measured the ulna and height of many women to see whether there is a
linear relationship between the two.

% TODO: the scatterplot from the study is not in the converted source. Re-capture it
% from the original worksheet and place it in xmPictures/linearModels/ --- problem 6
% below cannot be answered without it.

\begin{problem}
The researchers decided the data does have a linear relationship, and found the
trendline $y = 4.39x + 45.89$.

\begin{prompt}
\textbf{(a)} What does the $x$ value represent in this context?
\begin{freeResponse}
\end{freeResponse}

\begin{solution}
$x$ is the length of an individual's ulna, in cm.
\end{solution}
\end{prompt}

\begin{prompt}
\textbf{(b)} What does the $y$ value represent in this context?
\begin{freeResponse}
\end{freeResponse}

\begin{solution}
$y$ is the individual's height, in cm.
\end{solution}
\end{prompt}

\begin{prompt}
\textbf{(c)} What do we call the $4.39$ value in this function? What does it mean in
terms of the context?
\begin{freeResponse}
\end{freeResponse}

\begin{solution}
It is the slope, $\frac{4.39 \text{ cm of height}}{1 \text{ cm of ulna}}$. According
to the model, for every increase of 1 cm in ulna length, the model predicts an
increase of 4.39 cm in a woman's height.
\end{solution}
\end{prompt}

\begin{prompt}
\textbf{(d)} What do we call the $45.89$ value? What does it mean in this context?
\begin{freeResponse}
\end{freeResponse}

\begin{solution}
It is the $y$-intercept. In context: when the ulna length is 0 cm, the model predicts
a height of 45.89 cm. As before, this is outside any meaningful range of the data.
\end{solution}
\end{prompt}

\begin{prompt}
\textbf{(e)} The researchers specified that the resulting $y$-value from this
trendline could be off by $\pm 7.03$. What does this mean in terms of this context?
\begin{freeResponse}
\end{freeResponse}
\end{prompt}
\end{problem}

\begin{problem}
Pick a data point that might describe an individual with ``proportionally short
arms'' and circle it on the graph above. What are you looking for graphically, and why?

\begin{freeResponse}
\end{freeResponse}
\end{problem}

\end{document}
